\chapter{Software}
\label{sec:software}

All code is in our version of \banzai{} (branch \texttt{rm-improvements}; Appendix~\ref{sec:app_install}). The
list below gives, for each Python module of the package \texttt{banzai\_floyds}, what it does and where in this report
it is described; it is meant for readers who want to use or check the code. The diagnostics and archive tools need
only \textsc{numpy}, \textsc{scipy}, \textsc{astropy} and \textsc{matplotlib}.

\begin{description}
\item[\texttt{diagnostics.rm\_qa}] narrow-doublet and single-line fits, sky and telluric shifts, order-overlap ratio
(Chapter~\ref{sec:diagnostics}).
\item[\texttt{diagnostics.sky\_lines}] off-object sky spectrum from \texttt{BINNED2D}; adaptive top-hat centroids.
\item[\texttt{utils.arc\_refit\_utils}] arc refit: arc spectrum collapsed along the slit, line-group model (eqs.~\ref{eq:arcgroup}--\ref{eq:boxcar}),
correction (eq.~\ref{eq:arccorr}).
\item[\texttt{utils.wavelength\_offset\_utils}] telluric zero point (eqs.~\ref{eq:tell}, \ref{eq:tellcomb},
\ref{eq:transfer}).
\item[\texttt{wavelength\_offsets.ObjectWavelengthCorrector}] pipeline stage that applies the telluric zero point
of the object (runs after the extracted spectra are trimmed).
\item[\texttt{diagnostics.reextract}] archive products only (improved v1, Appendix~\ref{sec:app_v1}): undoing the
fringe division, master pattern, lamp-flat matching, harmonic fringe model, sky and extraction (eqs.~\ref{eq:undo}--\ref{eq:boxcarext});
command line \texttt{python -m banzai\_floyds.diagnostics.reextract}.
\item[\texttt{diagnostics.correct\_archive\_wavelengths}] arc refit and telluric zero point on archive
\texttt{-1d} files; console script \texttt{banzai\_floyds\_correct\_wavelengths}.
\item[\texttt{utils.sensitivity\_utils}, \texttt{diagnostics.recalibrate\_flux}] sensitivity as a template times a
nightly correction, and telluric correction, for archive products (eqs.~\ref{eq:sens}--\ref{eq:apply}).
\item[\texttt{utils.trace\_model\_utils}] object-trace model: refraction along the slit, instrumental template,
smooth correction of the residuals (eqs.~\ref{eq:adr}--\ref{eq:tracehybrid}); template in
\texttt{data/orders/trace\_template.dat}.
\item[\texttt{utils.harmonic\_fringe\_utils}] harmonic fringe model (shared by the pipeline and the archive tool),
including the blended phase (high-passed lamp phase in the blue part of the model, plain phase further red;
Section~\ref{sec:fringe_pipeline}).
\item[\texttt{fringe.HarmonicFringeCorrector}] pipeline stage that fits the harmonic fringe model; it runs after the
standards are loaded (stage \texttt{StandardLoader}). The stage \texttt{FringeLoader} supplies the phase from the
phase master.
\item[\texttt{telluric}] transmission applied only in the absorption bands; ensemble of the nearest standards for
science frames (\texttt{utils.telluric\_utils}).
\item[\texttt{flux}, \texttt{utils.flux\_utils}] ensemble flux calibration (grey-normalised median of the nearest
standards).
\item[\texttt{arc\_lines}] used arc lines, now including Hg\,4916 and Ar\,6416, 6677, 6753, 6767 with measured
strengths (\texttt{work/scripts/arc\_line\_strengths.py}).
\item[\texttt{data/standards}] reference spectra; Feige\,67 (extended beyond 9200\,\AA{} with a fitted power law)
and GD\,153 added.
\item[\texttt{offline}, \texttt{fringe\_masters}] offline reduction of a data set; construction of the phase masters
and of the residual fringe templates.
\end{description}
Analysis scripts used for this report (in \texttt{work/}): \texttt{scripts/v3\_verify.py} (step-by-step verification
of a reduction, Section~\ref{sec:sites_verify}), \texttt{scripts/compare\_versions.py} and
\texttt{make\_raw\_compare\_tex.py} (Table~\ref{tab:versions}), \texttt{scripts/build\_trace\_template.py} (trace
template and priors of a site), \texttt{scripts/arc\_check\_v3.py} (arc line-group shifts),
\texttt{v3\_chain.py} (a whole reduction in batches that keep the disk use low) and \texttt{banzai\_run/lean\_2d.py}
(lean 2-d products).
New settings (all defaults of this branch) choose the fringe method and its phase source, nodes, high-pass and
telluric blending; the object-trace model and its priors; and the telluric bands and the number and time span of the
standards in the ensemble: \texttt{FRINGE\_METHOD='harmonic'}, \texttt{FRINGE\_PHASE\_SOURCE='stacked'},
\texttt{FRINGE\_HARMONIC\_NODES}, \texttt{FRINGE\_PHASE\_HIGHPASS}, \texttt{FRINGE\_TELLURIC\_BLEND},
\texttt{PROFILE\_TRACE\_MODEL='adr'}, \texttt{PROFILE\_TRACE\_MODEL\_PRIORS}, \texttt{TELLURIC\_MODEL\_BANDS},
\texttt{TELLURIC\_ENSEMBLE\_N}, \texttt{TELLURIC\_ENSEMBLE\_DAYS}; setting them to the 1.1.0 values
(\texttt{'2d'}, \texttt{'polynomial'}, \texttt{None}) restores the stock behaviour. The offline reduction is the
command \texttt{banzai\_floyds\_offline} (subcommands \texttt{setup-db}, \texttt{skyflats}, \texttt{run},
\texttt{chain}; module \texttt{offline}), and the fringe products are built by \texttt{banzai\_floyds\_make\_phase\_master} and
\texttt{banzai\_floyds\_make\_residual\_template} (module \texttt{fringe\_masters}); Appendix~\ref{sec:app_install}.
From raw, all corrections run inside \banzai{}; \texttt{banzai\_floyds\_offline chain} reduces, in order, the arcs,
lamp flats, phase masters, standards, residual fringe template, the standards again, and the science frames. On
archive products (improved v1) the corrections are applied in this order: re-extraction $\rightarrow$ arc refit $\rightarrow$ telluric zero point $\rightarrow$ flux calibration. Unit tests with synthetic data (\texttt{tests/test\_wavelength\_offsets.py},
\texttt{test\_arc\_refit.py}, \texttt{test\_reextract\_flux.py}, \texttt{test\_trace\_model.py},
\texttt{test\_telluric.py}, \texttt{test\_flux\_ensemble.py}, \texttt{test\_fringe\_masters.py}) cover every new
step; the full suite of the package (176 tests) passes.
